\sectioncentered*{Заключение}
\addcontentsline{toc}{section}{Заключение}

Для регулярного выражения варианта~8 построен недетерминированный конечный автомат с $\varepsilon$-переходами, распознающий язык $a^{2}(ab)^{*}(c\mid d)^{+}$. Автомат недетерминирован из-за $\varepsilon$-дуг из состояний после префикса~$a^{2}$ и после тела цикла~$(ab)$.

Построением подмножеств получен детерминированный автомат из шести достижимых состояний. Он не минимален: два состояния, соответствующие позиции <<после $aa(ab)^{*}$>>, неотличимы и склеиваются. Минимальный ДКА содержит пять состояний. Восстановленное по нему регулярное выражение совпадает с исходным, что подтверждает корректность преобразований.

Поставленная цель достигнута: построены диаграмма и таблица исходного автомата, выполнены детерминизация и минимизация, получено эквивалентное регулярное выражение.

\begin{thebibliography}{9}
	\addcontentsline{toc}{section}{Список использованных источников}

	\bibitem{aho}
	Ахо А., Лам М., Сети Р., Ульман Дж. Компиляторы: принципы, технологии и инструментарий. \mdash{} 2-е изд. \mdash{} М. : Вильямс, 2008. \mdash{} 1184~с.

	\bibitem{hopcroft}
	Хопкрофт Д., Мотвани Р., Ульман Дж. Введение в теорию автоматов, языков и вычислений. \mdash{} 2-е изд. \mdash{} М. : Вильямс, 2002. \mdash{} 528~с.

	\bibitem{dragon-en}
	Aho A. V., Lam M. S., Sethi R., Ullman J. D. Compilers: Principles, Techniques, and Tools. \mdash{} 2nd ed. \mdash{} Boston : Addison-Wesley, 2006. \mdash{} 1000~p.

	\bibitem{yapis}
	Методические указания к практическим занятиям по дисциплине <<Языковые процессоры интеллектуальных систем>>. Задача~1. Построение конечного автомата.
\end{thebibliography}
